\chapter{紧致化}
\label{chap:compactifications-solutions}

\begin{solution}{local-compact-shrinking}
给定 \(x\in U\)，其中 \(U\) 是 \(X\) 中开集。
取 \(x\) 的紧邻域 \(K\) 和开集 \(W\)，使 \(x\in W\subseteq K\)。
\(W\cap U\) 是紧 Hausdorff 空间 \(K\) 中含 \(x\) 的相对开集。由紧
Hausdorff 空间的正则性，存在 \(K\) 中相对开集 \(O\) 使
\[
 x\in O\subseteq\overline O^{\,K}\subseteq W\cap U.
\]
写 \(O=K\cap G\)，其中 \(G\) 在 \(X\) 中开，并令 \(V=G\cap W\)。
由于 \(W\subseteq K\)，有 \(V\subseteq O\)，从而
\[
 \overline V^{\,X}\subseteq\overline O^{\,K}\subseteq U.
\]
而 \(\overline V^{\,X}\) 是紧集 \(K\) 的闭子集，故紧。
\end{solution}

\begin{solution}{locally-compact-open-subspace}
设 \(O\subseteq X\) 开，\(x\in O\)。由邻域收缩性质，存在 \(X\) 中开集
\(V\) 使
\[
 x\in V\subseteq\overline V^{\,X}\subseteq O,
 \qquad \overline V^{\,X}\text{ 紧}.
\]
此时 \(V\) 也是 \(O\) 中开集，并且
\(\overline V^{\,O}=\overline V^{\,X}\)，所以 \(x\) 在 \(O\) 中有紧邻域。

若 \(F\subseteq X\) 闭而 \(x\in F\)，取 \(X\) 中紧邻域 \(K\) 的内部
所含开集 \(W\ni x\)。则 \(K\cap F\) 是 \(F\) 中的紧邻域，因为它包含
相对开集 \(W\cap F\)。故闭子空间也局部紧。Hausdorff 性由子空间继承。
\end{solution}

\begin{solution}{one-point-neighborhoods}
若网 \(x_\alpha\to\infty\)，对每个紧集 \(K\subseteq X\)，集合
\(\{\infty\}\cup(X\setminus K)\) 是 \(\infty\) 的邻域，故网最终落入
其中；因为网取值于 \(X\)，这正是最终离开 \(K\)。反之，每个
\(\infty\) 的邻域都包含某个上述基本邻域，最终离开相应紧集便保证最终
进入该邻域。
\end{solution}

\begin{solution}{one-point-real-line-circle}
定义
\[
 \Phi(t)=\left(\frac{2t}{1+t^2},\frac{t^2-1}{1+t^2}\right)
 \quad(t\in\R),
 \qquad \Phi(\infty)=(0,1).
\]
直接计算知 \(\Phi(t)\in S^1\)。除北极外，逆映射为
\[
 t=\frac{x}{1-y},
\]
所以 \(\Phi:\R\to S^1\setminus\{(0,1)\}\) 是同胚。为检验新增点处
的连续性，直接计算
\[
 |\Phi(t)-(0,1)|=\frac{2}{\sqrt{1+t^2}}.
\]
给定 \(\varepsilon>0\)，取 \(R>2/\varepsilon\)。在 \(\infty\) 的
开邻域 \(\{\infty\}\cup(\R\setminus[-R,R])\) 上，像与北极的距离
小于 \(\varepsilon\)。故 \(\Phi\) 连续。它是紧空间到 Hausdorff 空间
的连续双射，因而是同胚。
\end{solution}

\begin{solution}{one-point-proper-extension}
先设 \(f\) 适当。给 \(\infty_Y\) 的基本邻域
\[
 W=\{\infty_Y\}\cup(Y\setminus K),
 \qquad K\subseteq Y\text{ 紧},
\]
其原像为
\[
 (f^*)^{-1}(W)
 =\{\infty_X\}\cup\bigl(X\setminus f^{-1}(K)\bigr).
\]
适当性保证 \(f^{-1}(K)\) 紧，所以这是 \(\infty_X\) 的邻域；在 \(X\)
内的连续性就是 \(f\) 的连续性。因此 \(f^*\) 连续。

反之，设 \(f^*\) 连续，\(K\subseteq Y\) 紧。集合
\(Y^*\setminus K\) 是 \(\infty_Y\) 的开邻域，故其原像是
\(\infty_X\) 的邻域。于是存在紧集 \(C\subseteq X\) 使
\[
 \{\infty_X\}\cup(X\setminus C)
 \subseteq (f^*)^{-1}(Y^*\setminus K).
\]
因而 \(f^{-1}(K)\subseteq C\)。它又是 \(X\) 中闭集，所以是紧集 \(C\)
的闭子集。故 \(f\) 适当。
\end{solution}

\begin{solution}{compactification-locally-compact-open}
若 \(X\) 在 \(K\) 中开，给 \(x\in X\)，利用紧 Hausdorff 空间的正则性取
\[
 x\in U\subseteq\overline U^{\,K}\subseteq X.
\]
则 \(\overline U^{\,K}\) 是 \(X\) 中的紧邻域，所以 \(X\) 局部紧。

反之，若 \(X\) 局部紧，给 \(x\in X\) 取 \(X\) 中开集 \(U\ni x\)，使
\(C=\overline U^{\,X}\) 紧。嵌入 \(X\hookrightarrow K\) 把 \(C\) 变成
\(K\) 中紧集，故 \(C\) 在 \(K\) 中闭。写 \(U=X\cap V\)，其中 \(V\)
在 \(K\) 中开。由于 \(X\cap V\subseteq C\)，开集 \(V\setminus C\)
与稠密集 \(X\) 不交，只能为空。因此
\[
 x\in V\subseteq C\subseteq X.
\]
每个 \(x\in X\) 都有一个 \(K\)-开邻域包含于 \(X\)，故 \(X\) 在 \(K\)
中开。
\end{solution}

\begin{solution}{stone-cech-bounded-extension}
若 \(|u|\le M\)，把 \(u\) 视为到紧区间 \([-M,M]\) 的连续映射。
Stone--Čech 泛性质给出连续延拓。对任意两个实值连续延拓 \(U_1,U_2\)，
等值集是连续映射 \((U_1,U_2):\beta X\to\R^2\) 对闭对角线的原像，
故在 \(\beta X\) 中闭。它含稠密集 \(j(X)\)，于是等于 \(\beta X\)。
这一唯一性论证不需预先假设候选延拓的值都落在 \([-M,M]\) 中。
\end{solution}

\begin{solution}{stone-cech-compact-fixed}
记嵌入为 \(j:X\to\beta X\)。因为 \(X\) 紧 Hausdorff，泛性质可用于
\(\operatorname{id}_X:X\to X\)，得到连续映射 \(r:\beta X\to X\)
满足 \(r\circ j=\operatorname{id}_X\)。映射 \(j\circ r\) 与
\(\operatorname{id}_{\beta X}\) 都是 \(j:X\to\beta X\) 的连续延拓，
由泛性质中的唯一性，\(j\circ r=\operatorname{id}_{\beta X}\)。因此
每个 \(z\in\beta X\) 都等于 \(j(r(z))\)，\(j\) 满射，且 \(r\) 是
它的连续逆映射。此证只用泛性质，不用评价像的具体构造。
\end{solution}

\begin{solution}{stone-cech-discrete-sequences}
离散空间上的任意函数连续。有界数列定义连续映射
\[
 a:\N\to[-M,M],\qquad a(n)=a_n,
\]
故唯一延拓到 \(\beta\N\)。

\(\N\) 的一点紧化要求函数在新增点连续，这等价于数列 \(a_n\) 在
\(n\to\infty\) 时有极限。一般有界数列如 \((-1)^n\) 没有极限，却仍能
延拓到 \(\beta\N\)，所以 Stone--Čech 的延拓性质严格更强。
\end{solution}

\begin{solution}{stone-cech-functoriality}
设 \(j_X:X\to\beta X\)、\(j_Y:Y\to\beta Y\) 是嵌入。复合
\[
 X\xrightarrow{f}Y\xrightarrow{j_Y}\beta Y
\]
以紧 Hausdorff 空间 \(\beta Y\) 为目标。由 \(\beta X\) 的泛性质，它
唯一延拓为 \(\beta f:\beta X\to\beta Y\)。

\(\beta(\operatorname{id}_X)\) 与 \(\operatorname{id}_{\beta X}\) 都延拓
\(j_X\)，故由唯一性相等。若 \(X\xrightarrow fY\xrightarrow gZ\)，则
\(\beta(g\circ f)\) 与 \(\beta g\circ\beta f\) 都把
\(j_X(x)\) 送到 \(j_Z(g(f(x)))\)，故二者也由唯一性相等。
\end{solution}

\begin{solution}{compactification-universal-uniqueness}
设 \((K_1,j_1)\)、\((K_2,j_2)\) 都满足泛性质。映射
\(j_2\circ j_1^{-1}:j_1(X)\to K_2\) 可视为从原空间 \(X\) 到 \(K_2\)
的连续映射，由 \(K_1\) 的泛性质延拓为 \(F:K_1\to K_2\)。同理得到
\(G:K_2\to K_1\)。复合 \(G\circ F\) 与
\(\operatorname{id}_{K_1}\) 在稠密集 \(j_1(X)\) 上相等，由延拓唯一性
处处相等；同理 \(F\circ G=\operatorname{id}_{K_2}\)。故 \(F\) 是保持
\(X\) 不变的同胚。

若 \(F_1,F_2:K_1\to K_2\) 都保持 \(X\) 不变，则它们在稠密集
\(j_1(X)\) 上相等。因为 \(K_2\) Hausdorff，等值集闭，所以
\(F_1=F_2\)。
\end{solution}
